\clearpage
\begin{center}
{\large SGA 7 --- EXPOS\'E XX}\par
\vspace{2em}
{\Large\underline{The Griffiths Theorem}}\par
\vspace{1.5em}
{\large by N. Katz}
\end{center}

\section*{0. Introduction}

We prove Griffiths's theorem~[1], which gives a systematic (though
nonconstructive) way of ``finding'' algebraic cycles that are
$\mathbb Q_\ell$-cohomologous to zero, but no multiple of which is
algebraically equivalent to zero.  The proof given here (due to
Grothendieck) is the translation into purely algebraic terms of Griffiths's
original, more or less transcendental, proof.  In both proofs, the
irreducibility theorem XVIII~6.7 is a key point (cf.~3.1.6 and~3.2).

\section*{1. The formalism of primitive classes}

Throughout what follows, a prime number $\ell$ is fixed.  All schemes under
consideration are assumed to have characteristic different from $\ell$.
Consider the situation

\begin{equation}\tag{1.0.1}
\begin{tikzcd}[column sep=huge, row sep=huge]
X \arrow[rr, "\pi"] \arrow[dr, "f"'] & {} & Y \arrow[dl, "g"] \\
{} & S & {}
\end{tikzcd}
\end{equation}
in which we assume
\begin{equation}\tag{1.0.3}
f,\ g \quad\text{smooth of relative dimensions}\quad d_1\ \text{and}\ d_2,
\end{equation}
and
\begin{equation}\tag{1.0.4}
\pi\quad\text{proper}.
\end{equation}
Since $X$ and $Y$ are smooth over $S$, the properness of $\pi$ gives a
\underline{Gysin morphism} (SGA~5, Expos\'e~IV; the further citation is left
blank in the print)
\begin{equation}\tag{1.0.5}
\pi_*:H^q(X,\mathbb Q_\ell(i))\longrightarrow
H^{q-2(d_1-d_2)}(Y,\mathbb Q_\ell(i-d_1+d_2)).
\end{equation}
Localizing on $S$ in the \`etale topology, we likewise obtain a morphism of
sheaves on $S$\footnote{In (1.0.6) the print has an apostrophe in place of
the closing parenthesis in the exponent and has $i-d_1-d_2$ in the target
twist.  Formulae (1.0.5) and (1.0.7), as well as the Gysin degree, give the
reading displayed here.}
\begin{equation}\tag{1.0.6}
\pi_*:R^qf_*\mathbb Q_\ell(i)\longrightarrow
R^{q-2(d_1-d_2)}g_*\mathbb Q_\ell(i-d_1+d_2),
\end{equation}
and, on applying the functor $H^p(S,-)$, a morphism
\begin{equation}\tag{1.0.7}
\pi_*:H^p(S,R^qf_*\mathbb Q_\ell(i))\longrightarrow
H^p(S,R^{q-2(d_1-d_2)}g_*\mathbb Q_\ell(i-d_1+d_2)).
\end{equation}

\subsubsection*{Definition 1.0.8.}
Let $(E_r^{p,q},d_r^{p,q})$ and
$({}'E_r^{p,q},{}'d_r^{p,q})$ be two spectral sequences in a given abelian
category, and let $(a,b)\in\mathbb Z\times\mathbb Z$.  The shift of
$({}'E_r^{p,q},{}'d_r^{p,q})$ by $(a,b)$ is the spectral sequence
$({}''E_r^{p,q},{}''d_r^{p,q})$ given by\footnote{The print omits the prime
on the right-hand differential in (1.0.9).  The definition of the shift of
the primed spectral sequence requires the primed differential.}
\begin{equation}\tag{1.0.9}
\left\{
\begin{aligned}
{}''E_r^{p,q}&={}'E_r^{p+a,q+b},\\
{}''d_r^{p,q}&={}'d_r^{p+a,q+b}.
\end{aligned}
\right.
\end{equation}

\subsubsection*{1.0.10.}
A morphism of spectral sequences from
$(E_r^{p,q},d_r^{p,q})$ to $({}''E_r^{p,q},{}''d_r^{p,q})$ is called a
morphism of bidegree $(a,b)$ from $(E_r^{p,q},d_r^{p,q})$ to
$({}'E_r^{p,q},{}'d_r^{p,q})$.

The following lemma is an immediate consequence of the theory of the
\underline{Gysin morphism} (SGA~5, Expos\'e~IV).

\subsubsection*{Lemma 1.1.}
There is a canonical morphism $\pi_*$ of bidegree
$(0,-2(d_1-d_2))$\footnote{The print has $(a,-2(d_1-d_2))$.  The map
(1.0.7) leaves the first Leray index $p$ unchanged, so the first component
of the bidegree is $0$.} between the Leray spectral sequences
\begin{equation}\tag{1.1.0}
R^p\alpha_*R^qf_*\mathbb Q_\ell(i)\Longrightarrow
R^{p+q}(\alpha f)_*\mathbb Q_\ell(i)
\end{equation}
and
\begin{equation}\tag{1.1.1}
R^p\alpha_*R^qg_*\mathbb Q_\ell(i-d_1+d_2)\Longrightarrow
R^{p+q}(\alpha g)_*\mathbb Q_\ell(i-d_1+d_2),
\end{equation}
which is given on the $E_2$ term by (1.0.7), and on the abutment by
(1.0.5).

\subsubsection*{1.2.}
Now consider the situation
\begin{equation}\tag{1.2.1}
\begin{tikzcd}[column sep=large, row sep=large]
{} & X\times_S Y \arrow[dl, "\mathrm{pr}_1"']
                  \arrow[dr, "\mathrm{pr}_2"] & {} \\
X \arrow[dr, "f"'] & {} & Y \arrow[dl, "g"] \\
{} & S \arrow[d, "\alpha"] & {} \\
{} & T & {}
\end{tikzcd}
\end{equation}
in which $S,X,f,Y,g$ are as in (1.0.1)--(1.0.4), but now $f$ is also
assumed proper.\footnote{The print begins this range at the nonexistent
(1.0.2); the displayed situation is (1.0.1).}  Let
\begin{equation}\tag{1.2.2}
z\in H^a(X\times_S Y,\mathbb Q_\ell(b)).
\end{equation}
Then $z$ induces a map, again denoted by $z$,
\begin{equation}\tag{1.2.3}
z:H^q(X,\mathbb Q_\ell(i))\longrightarrow
H^{q+a-2d_1}(Y,\mathbb Q_\ell(i+b-d_1))
\end{equation}
by the formula
\begin{equation}\tag{1.2.4}
z(x)=\mathrm{pr}_{2*}\bigl(z\wedge\mathrm{pr}_1^*(x)\bigr),
\end{equation}
where $\mathrm{pr}_{2*}$ is the Gysin homomorphism.  By \`etale localization
on $S$, the class $z$ likewise defines
\begin{equation}\tag{1.2.5}
z:R^qf_*\mathbb Q_\ell(i)\longrightarrow
R^{q+a-2d_1}g_*\mathbb Q_\ell(i+b-d_1),
\end{equation}
and, after applying $H^p(S,-)$, maps
\begin{equation}\tag{1.2.6}
H^p(S,R^qf_*\mathbb Q_\ell(i))\longrightarrow
H^p(S,R^{q+a-2d_1}g_*\mathbb Q_\ell(i+b-d_1)).
\end{equation}

\subsubsection*{Lemma 1.3.}
The class $z$ defines a morphism of bidegree
$(0,a-2d_1)$\footnote{The print has $(0,b-2d_1)$.  The maps (1.2.3),
(1.2.5), and (1.2.6) change cohomological degree by $a-2d_1$; $b$ controls
the Tate twist.} between the Leray spectral sequences
\begin{equation}\tag{1.2.7}
H^p(S,R^qf_*\mathbb Q_\ell(i))\Longrightarrow
H^{p+q}(X,\mathbb Q_\ell(i))
\end{equation}
and
\begin{equation}\tag{1.2.8}
H^p(S,R^qg_*\mathbb Q_\ell(i+b-d_1))\Longrightarrow
H^{p+q}(Y,\mathbb Q_\ell(i+b-d_1)),
\end{equation}
which is given on the $E_2$ terms by (1.2.6), and on the abutments by
(1.2.3).

\noindent\underline{Proof.}
In view of (1.2.4), and since ``multiplication by $z$'' defines a morphism
of bidegree $(0,a)$ between the Leray spectral sequences
\begin{equation}\tag{1.2.9}
H^p\bigl(S,R^q(f\circ\mathrm{pr}_1)_*\mathbb Q_\ell(i)\bigr)
\Longrightarrow H^{p+q}(X\times_S Y,\mathbb Q_\ell(i))
\end{equation}
\begin{center}
\begin{tikzcd}
{} \arrow[d, "z"'] \\
{}
\end{tikzcd}
\end{center}
\begin{equation}\tag{1.2.10}
H^p\bigl(S,R^{q+a}(f\circ\mathrm{pr}_1)_*\mathbb Q_\ell(i+b)\bigr)
\Longrightarrow H^{p+q+a}(X\times_S Y,\mathbb Q_\ell(i+b)),
\end{equation}
it suffices to apply Lemma~1.1 to the situation\footnote{In (1.2.9) and
(1.2.10) the print writes $(\mathrm{pr}_1 f)_*$, whereas the map to $S$ is
$f\circ\mathrm{pr}_1$, as the printed diagram (1.2.11) itself shows.}
\begin{equation}\tag{1.2.11}
\begin{tikzcd}[column sep=large, row sep=large]
X\times_S Y \arrow[rr, "\mathrm{pr}_2"]
             \arrow[dr, "f\circ\mathrm{pr}_1"]
  & {} & Y \arrow[dl, "g"] \\
{} & S & {}
\end{tikzcd}
\qquad .
\end{equation}

\subsubsection*{Corollary 1.4.}
In the situation (1.2.1)--(1.2.2), the morphism (1.2.3)
\begin{equation}\tag{1.4.1}
z:H^q(X,\mathbb Q_\ell(i))\longrightarrow
H^{q+a-2d_1}(Y,\mathbb Q_\ell(i+b-d_1))
\end{equation}
takes primitive classes (XVIII~5.8.1) to primitive classes, and the following
diagram is commutative:
\begin{equation}\tag{1.4.2}
\begin{tikzcd}[column sep=large, row sep=large]
\mathrm{Prim}^q(X/S,\mathbb Q_\ell(i))
  \arrow[r, "z"] \arrow[d, "u_f"]
& \mathrm{Prim}^{q+a-2d_1}(Y/S,\mathbb Q_\ell(i+b-d_1))
  \arrow[d, "u_g"] \\
H^1(S,R^{q-1}f_*\mathbb Q_\ell(i))
  \arrow[r, "z"]
& H^1(S,R^{q+a-2d_1-1}g_*\mathbb Q_\ell(i+b-d_1)).
\end{tikzcd}
\end{equation}
Here $u_f$ and $u_g$ are the canonical homomorphisms of XVIII~5.8.1.
